% Chapter 15: Training Regimes, Optimization, and Failure Modes
% Part II of the book.

\chapter{Training Regimes, Optimization, and Failure Modes}
\label{ch:part2-training-regimes-optimization-and-failure-}
\label{sec:training}
% -----------------------------------------------------------------------------

Part~I's error budget stops at the edge of arithmetic. Theorem~\ref{thm:fno_convergence}
bounds the distance between a target solution operator and \emph{some} member of the
Fourier neural operator class; it says nothing about the distance between that member and
the parameter vector a training loop hands back, and \eqref{eq:ch05-error-decomposition}
deliberately contains no term for the optimizer, the numerical format, or the batch order.
That gap, between the best parameters a class contains and the ones a loop returns, is where
surrogates are lost in practice, and it has three faces. A spectral layer's parameters carry
gradients whose magnitude falls algebraically across the array, so a step size is never one
number but a profile. The coefficient array a transform produces spans a dynamic range that
collides with what half-precision arithmetic can store. And a network accurate on one
application of the time-$\Delta t$ map can be worthless after fifty, because
\eqref{eq:ch09-gronwall} converts a one-step error into a trajectory error at a rate set by
$\Lambda_\theta$, which \eqref{eq:ch09-lipschitz} makes an explicit functional of the weights
and no training objective in this book controls.

This chapter takes the three up in turn and collects the residue as a field guide.
Learning-rate scheduling comes first, because the schedule conditions the multiplier array
and because a cyclic cosine anneal interacts with the mode budget in a way that changes how
a mode-count experiment should be read. Mixed precision follows, with the arithmetic of the
two half-precision formats set out explicitly, since the question worth asking is not whether
running in reduced precision is faster but which representable range a spectral accumulation
requires. Autoregressive rollout and pushforward stabilization come third, as the engineering
counterpart of Part~I's
unrolled error analysis: a recipe, not a theorem, and used here as such.
\index{training regimes!schedule, format, and horizon}

The reader should bring PyTorch at the level of \cite{paszke2019pytorch} --- tensors,
modules, an autograd backward pass --- and the vocabulary of a graphics processing unit
(GPU), of CUDA (Compute Unified Device Architecture), its programming platform, and of cuDNN
(the CUDA Deep Neural Network library), which services the pointwise convolutions used for
lifting and projection. Two objects from earlier in this part are used without
re-derivation: the relative $L^2$ loss \eqref{eq:rel_l2_loss} of
\chapref{ch:part2-the-practitioner-s-paradigm-infinite-dimen}, and the
\texttt{SpectralConv2d} layer of \chapref{ch:part2-fourier-neural-operators-fno-from-scratch},
whose paired complex weight tensors are the multipliers the first section below
differentiates. From Part~I come the multiplier of \eqref{eq:ch05-multiplier}, the retained
set $Z_{k_{\max}}$ of \eqref{eq:ch05-modeset}, and the coefficient decay
\eqref{eq:ch05-coefficient-decay}. Recipes are consolidated in Appendix~\ref{app:d}; nothing
in this chapter is a measured outcome, and where a mechanism is claimed it is arithmetic over
a library convention, said as such.

% -----------------------------------------------------------------------------
\section{Learning Rate Scheduling in Spectral Space}
\label{sec:ch26-schedule}
% -----------------------------------------------------------------------------

Why a spectral layer needs a schedule that an ordinary convolutional network can do
without is visible in one line of algebra. Let $v_\ell$ be the input of layer $\ell$,
$\hat{v}_\ell(k)$ its coefficients on the retained set $Z_{k_{\max}}$,
$R_{\phi,\ell}(k)$ the complex multiplier of \eqref{eq:ch05-multiplier}, and
$\hat{\lambda}_\ell(k)$ the coefficients of the adjoint field
$\partial\mathcal{L}/\partial v_{\ell+1}$ that the backward pass delivers in physical
space. Because the layer acts mode by mode, the derivative of the loss with respect to the
multiplier at a single mode factorizes:
\begin{equation}
  \label{eq:ch26-mode-gradient}
  \frac{\partial \mathcal{L}}{\partial R_{\phi,\ell}(k)} \;=\; \hat{\lambda}_\ell(k)\, \hat{v}_\ell(k)^{\mathsf{T}}, \qquad k \in Z_{k_{\max}}.
\end{equation}
Two qualifications accompany \eqref{eq:ch26-mode-gradient}. It holds up to the constant the
transform convention attaches to the adjoint, and up to the treatment of a complex tensor as
a pair of real ones, which is how the framework of \cite{paszke2019pytorch} resolves the
derivative. Its norm is bounded by the product $\lvert\hat{v}_\ell(k)\rvert\,
\lvert\hat{\lambda}_\ell(k)\rvert$, and for the smooth inputs the convergence theory
assumes, both factors decay algebraically in $\lvert k\rvert$ by
\eqref{eq:ch05-coefficient-decay}. Gradient magnitude therefore spans orders of magnitude
across one parameter array, and it is smallest at the band edge --- precisely at the modes
whose removal the truncation term $C_1 k_{\max}^{-\alpha}$ of
\eqref{eq:ch05-error-decomposition} charges to the model.
\index{Fourier multiplier $R_\phi(k)$!gradient of}
\index{condition number!across the mode budget}

The factorization is a property of the layer, not of the network, and the distinction
matters. The activation acts pointwise in physical space, which scatters energy among
frequencies, so a gradient entering at one mode returns at all of them;
\eqref{eq:ch26-mode-gradient} gives the layer-local scaling of the parameter-to-loss map,
not a decoupled optimization problem. The scaling is also the reason a constant step size is
the wrong instrument here, where a conventional network tolerates one: a single increment
applied to an array whose usable magnitudes differ by orders moves the low modes a useful
distance while displacing the band-edge multipliers by an amount large relative to their own
gradient signal, and early training destabilizes. \index{learning rate scheduling}

Adam (adaptive moment estimation) is what makes such an array trainable at all: each
coordinate is advanced by its gradient divided by a running root-mean-square of recent
gradients, so every mode moves at a comparable absolute rate regardless of its own signal
strength. The benefit carries a cost that is easy to forget. The same normalization divides
away the \emph{reliability} of the gradient, so band-edge parameters whose gradients sit at
the level of minibatch noise are stepped at full rate on noise. AdamW retains that
normalization of the second moment but applies weight decay to the parameters directly rather
than through it. On a complex multiplier, decay acts on the real and imaginary components
jointly: it shrinks the magnitude of $R_{\phi,\ell}(k)$ and leaves its phase untouched.
Through
\eqref{eq:ch09-lipschitz}, and the $L^2$ version recorded in Remark~\ref{rem:ch05-lipschitz},
decay is thus the one mechanism by which the training loop touches the Lipschitz product
Part~I said no loss controls --- and one more way in which the band edge is penalized,
since decay shrinks every mode toward zero at the same rate whether or not that mode
carries signal. \index{AdamW}\index{weight decay!on complex multipliers}

Listing~\ref{lst:ch26-onecycle} is the reference configuration, and the reader should take
four decisions out of it rather than the code itself. First, the peak rate
\texttt{max\_lr} is the only absolute learning rate the scheduler is given: the cycle starts
at one twenty-fifth of the peak, $2\times10^{-4}$, rises to $5\times10^{-3}$ over the first
fifth of the run (\texttt{pct\_start=0.2}), and anneals to one thousandth of the starting
value, $2\times10^{-7}$; the \texttt{lr=1e-3} passed to \texttt{AdamW} is overwritten on
the first scheduler call and should be read as a placeholder. Second, \texttt{total\_steps}
is counted in iterations, and \texttt{scheduler.step()} is placed inside the batch loop, so
the two must agree --- mismatch them and the cycle completes early, after which the rate
sits at its floor for the rest of the run while the curve looks like a plateau. Third,
clipping to a global gradient norm of \texttt{max\_norm=1.0} is computed over all
parameters jointly, so a single band-edge mode with an outlier gradient shrinks the update
of every other tensor in the model; and because Adam normalizes by its own second moment, a
clip factor that persists across steps is largely absorbed, which is why clipping buys
protection against transients rather than against a systematically ill-scaled array.
Fourth, the default annealing law of \texttt{OneCycleLR} is a cosine, so the cyclic schedule
\emph{is} a cosine decay wrapped in a short warmup --- which is where the interaction with
truncation becomes visible. \index{OneCycleLR}\index{gradient clipping}

Annealing and the mode budget pull in opposite directions. Late in a run, the content still
available to be fitted is the small-amplitude content: the low modes, the bypass $W_\ell$,
and the lifting and projection maps have converged against their own signal several epochs
earlier, and what remains is the band edge, where \eqref{eq:ch26-mode-gradient} says the
gradients are smallest. A schedule whose rate falls by four orders of magnitude across the
remaining four fifths of the run suppresses exactly those parameters while the loss
reduction still available is coming from them. The practical consequence is that $k_{\max}$
and the schedule shape are one decision, not two: a wider band must be given proportionally
more optimizer steps before its annealed tail counts as a fair trial. Doubling the retained band while holding \texttt{total\_steps}
and \texttt{final\_div\_factor} fixed adds parameters that receive their share of
optimization time from the annealed tail; a run whose validation error stops improving as
the rate decays has not necessarily reached the truncation floor of
Figure~\ref{fig:spectral_decay}, and may simply have run out of learning rate at the band
edge. Section~\ref{sec:ch26-failures} gives the test that separates the two readings.
\index{learning rate scheduling!interaction with mode truncation}
\index{cosine annealing}

One further coupling deserves recording because it defeats practitioners who port a recipe
between grids. The transform called in
\chapref{ch:part2-fourier-neural-operators-fno-from-scratch} uses the library default
normalization, under which the forward direction is unscaled and the inverse divides by the
number of grid points; coefficients of a fixed continuous field sampled on a finer grid
therefore grow in absolute magnitude with the grid, and so do the products in
\eqref{eq:ch26-mode-gradient}. The parameter shapes and the architecture are untouched by
refinement --- that invariance is the subject of
\chapref{ch:part1-discretization-invariance-and-continuum-li} --- but the ratio between
per-mode gradient magnitude and minibatch noise is not, so a peak rate and a decay
coefficient tuned at one resolution require rechecking at another. Requesting orthonormal
transform normalization removes the grid factor from the coefficient array, because the
library then divides by the square root of the number of points in each direction. That
change does not reparameterize the network: the two factors cancel in the round trip, as
they must for an identity multiplier, so the same $R_\phi$ computes the same map. What it
does change is the gradient, since both factors in \eqref{eq:ch26-mode-gradient} shrink by
the square root of the grid and their product by the whole grid. Adam absorbs a uniform
rescaling of the gradient --- its update is a ratio of a moment to the root mean square of
the same quantity --- which is why the switch passes unnoticed most of the time; a fixed
\texttt{weight\_decay} does not, because decay is written against the parameters rather than
against the gradient, so shrinking the gradient by the grid size silently strengthens decay's
grip on exactly the band edge that has the least signal to begin with. The next section makes
the same point from the storage side, where the grid factor is not a tuning matter but an
overflow one.

\begin{lstlisting}[language=Python, label={lst:ch26-onecycle}, caption={Reference optimizer and schedule configuration for a spectral network: decoupled weight decay, a cyclic schedule with cosine annealing whose initial, peak and terminal learning rates are fixed by the two divisor arguments, and a global gradient-norm clip applied before every update.}]
optimizer = torch.optim.AdamW(model.parameters(), lr=1e-3, weight_decay=1e-4)
scheduler = torch.optim.lr_scheduler.OneCycleLR(
    optimizer, max_lr=5e-3, total_steps=epochs * len(train_loader),
    pct_start=0.2, div_factor=25.0, final_div_factor=1000.0
)

# Training loop step
for x, y in train_loader:
    optimizer.zero_grad()
    pred = model(x)
    loss = loss_fn(pred, y)
    loss.backward()
    torch.nn.utils.clip_grad_norm_(model.parameters(), max_norm=1.0)
    optimizer.step()
    scheduler.step()
\end{lstlisting}

% -----------------------------------------------------------------------------
\section{Mixed Precision (AMP) Considerations}
\label{sec:ch26-amp}
% -----------------------------------------------------------------------------

Automatic mixed precision (AMP) executes most of a network's arithmetic in a 16-bit
floating-point format while holding a single-precision copy of the parameters; on the GPUs
this part targets it is the standard route to larger batches. For operator models the
question is not whether to use it but where its two candidate formats stop short of the
arithmetic a spectral layer performs, and the axis the answer turns on is representable range
against precision: the formats trade those away differently, and a Fourier coefficient array
stresses both ends at once. Table~\ref{tab:ch26-formats} gives the constants; the conclusion
to take from it is that \texttt{bfloat16} differs from \texttt{float32} in the significand
alone, whereas \texttt{float16} differs in exponent and significand both, so half precision
fails for reasons no tuning of the model can remove. \index{automatic mixed precision (AMP)}
\index{floating-point formats!dynamic range versus precision}

\begin{table}[htbp]
\centering
\small
\caption{The three formats a spectral layer touches. Bit counts and powers of two are exact;
the largest finite values are rounded to three figures. \texttt{bfloat16} is the 16-bit
truncation of \texttt{float32} --- same exponent width, sixteen fewer significand bits. Unit
roundoff is half the spacing of representable numbers near one.}
\label{tab:ch26-formats}
\begin{tabular}{lccc}
\toprule
& \textbf{\texttt{float16}} & \textbf{\texttt{bfloat16}} & \textbf{\texttt{float32}} \\
\midrule
Exponent bits & $5$ & $8$ & $8$ \\
Explicit significand bits & $10$ & $7$ & $23$ \\
Unit roundoff & $2^{-11}$ & $2^{-8}$ & $2^{-24}$ \\
Largest finite value & $65\,504$ & $\approx 3.39\times 10^{38}$ & $\approx 3.40\times 10^{38}$ \\
Least normal value & $2^{-14}\approx 6.1\times 10^{-5}$ & $2^{-126}\approx 1.2\times 10^{-38}$ & $2^{-126}\approx 1.2\times 10^{-38}$ \\
Bytes per real array entry & $2$ & $2$ & $4$ \\
\bottomrule
\end{tabular}
\end{table}

Begin with the complex weights, where the format question has a clean answer. The multiplier
tensors are allocated in \texttt{torch.cfloat}, two single-precision numbers per entry, and it
is tempting to read AMP as halving that storage: two bytes per component instead of four, so
the dense complex output array of shape $(B, C_{\mathrm{out}}, H, \lfloor W/2 \rfloor + 1)$
that the layer fills on every call would cost half as much. The reading is wrong in both
directions at once. AMP keeps the master parameters and Adam's moments in single precision by
construction, so a policy that demotes arithmetic never touches the stored multiplier: nothing
there is reclaimable, and equally nothing there is lost. What the policy reaches is the
arithmetic inside the call, and there precision and range must be argued separately. Take
precision first. The per-mode contraction --- an \texttt{einsum}, a sum in Einstein summation
convention over the channel
axis of the frequency slice --- accumulates $C_{\mathrm{in}}$ addends, each a product of a
coefficient with a weight, and in a binary16 accumulator the unit roundoff is $2^{-11}$: once
the running partial sum exceeds an addend by more than about two thousand to one, the addend
is discarded by the arithmetic rather than by the model. Were the contraction carried out in
half precision, each component of $R_{\phi,\ell}(k)$ would additionally keep eleven
significant bits, which bounds how finely the trained filter can place a notch or a phase
shift --- the quantities that distinguish a learned Green's function from a smooth blur.
Range is the other end of the same format, and it is the larger of the two problems.
\index{complex weights!half precision}
\index{einsum (Einstein summation contraction)}

The intermediate arrays are where the collision actually occurs, and the collision is
arithmetic, not anecdotal. \texttt{torch.fft.rfft2} denotes the two-dimensional
real-to-half-complex fast Fourier transform (FFT), and with the library default the forward
transform is unscaled: the
zero-frequency, or direct-current (DC), coefficient of a field is the sum of all its values.
On a $256\times256$ grid, a field of unit spatial mean therefore produces a DC coefficient of
$65\,536$, which exceeds the largest representable \texttt{float16} value, $65\,504$, by
thirty-two. Normalizing the ensemble pointwise as
\chapref{ch:part2-the-practitioner-s-paradigm-infinite-dimen} prescribes does not repair
this, because ensemble normalization removes the mean \emph{at each location across samples}
and leaves each sample's own spatial mean standing. Refine the grid and the whole coefficient
array shifts upward: the same continuous field sampled four times as finely in each
direction multiplies every fixed-physical-wavenumber coefficient by roughly sixteen. The
other end of the same array fails by a different arithmetic: one storage must hold the zero
mode and the band edge together, and the usable span of binary16 is about $2^{40}$ counting
subnormals, while the ratio between the two is set not by the grid but by the field's
spectrum. For a decay as steep as \eqref{eq:ch05-coefficient-decay} allows, coefficients near
$k_{\max}$ fall below the least normal value $2^{-14}$ while the DC term sits comfortably
under the ceiling. Two independent quantities --- grid size and spectral slope --- therefore
decide which end breaks, and neither is a property of the model. That is
the mechanism behind the warning attached to \texttt{torch.fft.rfft2} in half precision, that
high-dynamic-range coefficient arrays under- and overflow; and it is why the failure appears
with a refined training grid on a recipe that ran at the coarse one. The inverse transform
divides by the grid size, so the physical output is back in range while the frequency
intermediate was not --- and it is why a loss can turn into a non-number (\texttt{NaN}) at
the second spectral layer while the pointwise lift that feeds it stays finite.
\index{float16!overflow in spectral sums}
\index{underflow!flush-to-zero at the band edge}
\index{rfft2 (two-dimensional real FFT)}

\texttt{bfloat16} removes the exponent half of the problem and none of the mantissa: eight
exponent bits give it the range of \texttt{float32}, so no accumulation of coefficients can
overflow or underflow that a single-precision run would survive, while the unit roundoff
$2^{-8}$ is three bits coarser than binary16's. The two costs are not symmetric. A radix-2
transform of $N$ points runs $\log_2 N$ butterfly stages and each stage rounds, so a narrow
significand is charged once per stage and the error it leaves behind grows like $\log_2 N$
times the unit roundoff: mild, bounded, and slow-growing in $N$. A range that overflows is
charged once and irrecoverably --- one coefficient leaves the representable set, and every
operation downstream inherits the infinity or the flush to zero. Note what the asymmetry does
not say. It does not make \texttt{bfloat16} free inside the layer, because that per-stage
charge is levied inside the transform, on quantities every later stage reads, and a
two-dimensional transform of a $256\times256$ grid is $\log_2 2^{16} = 16$ such roundings of
$2^{-8}$ --- a worst-case relative error of $6\times10^{-2}$, several times the error floors
this part reports. This is the reasoning behind the recommended policy for operator
models: run the lifting and projection convolutions, which are pointwise $1\times1$
maps with no global accumulation, in \texttt{bfloat16}, and keep the transform itself and the
complex multiplication in \texttt{float32}. Two details of that policy are worth more than
the policy. It must be expressed as an explicit dtype cast at the boundary of the spectral
layer, or by disabling autocast around that layer, because
\texttt{torch.cuda.amp.autocast} assigns dtypes by operation type rather than by the role an
operation plays in the model: the per-mode contraction reaches it as a sum over channels,
indistinguishable there from any other batched product, and no operation table records that
this particular sum is the one whose dynamic range was just measured. And
because the spectral and bypass paths of the layer meet in an addition, casting a demoted
bypass back up to \texttt{float32} at that point recovers nothing, since the precision of a
sum is set by its coarsest addend. The savings are real where nothing accumulates --- the
$1\times1$ convolutions, the local map $W_\ell$, the stored activations --- and they are
illusory wherever the mode contraction reaches. \index{bfloat16}\index{CUDA}
\index{cuDNN (CUDA Deep Neural Network library)}

Three situations lie outside what AMP can be configured around. The first is any training
functional that evaluates a PDE (partial differential equation) residual by finite
differencing or automatic differentiation of the output field, as in the physics-informed
neural operator (PINO) objective \eqref{eq:pino_loss}: differencing subtracts nearly equal
numbers, so a field stored with $2^{-8}$ relative resolution yields a derivative whose error
is set by the storage format rather than by the model, and a wider exponent does nothing about
cancellation. The second is the
metric. The relative $L^2$ loss \eqref{eq:rel_l2_loss} is a ratio of two reductions over
$B\times C\times H\times W$ entries, and accumulated in a half-precision format it carries a
rounding error no smaller than the unit roundoff of Table~\ref{tab:ch26-formats} ---
$2^{-8} = 3.9\times10^{-3}$ for \texttt{bfloat16} --- which is of the same order as the
percent-level error floors \chapref{ch:part2-canonical-pde-benchmarks-and-reference-wor}
reports. Model error and arithmetic then become indistinguishable; losses, and every reported
error, belong in single precision. The third is binary16 gradient range: the backward pass
produces values orders of magnitude below the forward ones, and the fix is loss scaling ---
multiply the loss by a large power of two, step, and unscale --- rather than any change to the
model; \texttt{bfloat16} needs no scaler precisely because its exponent already matches
\texttt{float32}. Appendix~\ref{app:d} records which of these settings each benchmark recipe
uses. \index{loss scaling (gradient scaling)}

% -----------------------------------------------------------------------------
\section{Autoregressive Rollout and Pushforward Stabilization}
\label{sec:ch26-rollout}
% -----------------------------------------------------------------------------

For an evolutionary problem the trained object is the time-$\Delta t$ map, applied to its
own output: with $v_0 = u_0$ and the network $\mathcal{G}_\theta$ standing in for one step
of the solver,
\begin{equation}
  \label{eq:ch26-rollout}
  v_{k+1} \;=\; \mathcal{G}_\theta(v_k), \qquad t_k = k\,\Delta t, \qquad k = 0, \dots, K-1,
\end{equation}
so that the $k$-th term of \eqref{eq:ch26-pushforward} below is $\mathcal{G}_\theta^k(u_0)
= v_k$. Part~I priced this recursion in \eqref{eq:ch09-onesided} and \eqref{eq:ch09-gronwall}
in terms of the one-step error $\tau$ at the true states and the global Lipschitz constant
$\Lambda_\theta$, and the arithmetic is unforgiving: if $\Lambda_\theta$ exceeds one, the
error bound grows geometrically in $K$ no matter how small $\tau$ is. The reference
configuration of \chapref{ch:part2-canonical-pde-benchmarks-and-reference-wor} unrolls the
vorticity form of the two-dimensional incompressible Navier--Stokes equations
\eqref{eq:ch03-navierstokes} across fifty to one hundred steps, and the global weather
surrogate of \cite{pathak2022fourcastnet} advances six-hour steps over days, so this is not
a corner of the training loop; it is the regime in which a time-marching operator is
actually used.
\index{autoregressive rollout!training on the model's own iterates}

The naive procedure is teacher forcing: sample pairs $\big(u(t_k), u(t_{k+1})\big)$ from
solver trajectories and minimize a one-step loss. Its defect is a mismatch of input
distributions rather than of arithmetic. Under \eqref{eq:ch26-rollout} at inference the
network is queried at $v_k$, which lies at distance $e_k$ from $u(t_k)$ --- the pathwise
error of \eqref{eq:ch09-onesided}, and generally not small once $k$ is past the first few
steps; the loss was minimized over a neighborhood of the data manifold, and
nothing constrains $\mathcal{G}_\theta$ on points that neighborhood excludes. Exposure bias
is the name for the resulting drift, and the cure implemented in practice is to train on the
model's own iterates --- the \textbf{pushforward} objective, which adds to the one-step loss
the weighted sum over unrolled powers
\begin{equation}
  \label{eq:ch26-pushforward}
  \mathcal{L}_{\text{pushforward}} = \sum_{k=1}^K \lambda_k \big\| \mathcal{G}_\theta^k(u_0) - u(k\Delta t) \big\|_{L^2}^2,
\end{equation}
so that the network is required to remain accurate, and by extension stable, on the
trajectories it actually generates. Equation \eqref{eq:ch26-pushforward} is written in the
$L^2$-squared form because that is what the automatic differentiation machinery handles
without a division by a field norm; the relative error a reader reports should still come
from \eqref{eq:rel_l2_loss}, and the two are not interchangeable at the level of which
samples dominate the gradient (Section~\ref{sec:ch26-failures}).
\index{teacher forcing}\index{exposure bias}
\index{pushforward regularization!weights $\lambda_k$ and horizon $K$}

That pushforward works is an empirical claim, and the mechanism is a change of sampling distribution rather than a new
regularizer. The quantity $\tau$ in
\eqref{eq:ch09-onesided} is a supremum over the true trajectory;
\eqref{eq:ch09-gronwall} then propagates it. Pushforward terms do not bound
$\Lambda_\theta$ --- no term in \eqref{eq:ch26-pushforward} contains a derivative of
$\mathcal{G}_\theta$ --- but they do place training inputs at the states the rollout
visits, which is where finite differences of $\mathcal{G}_\theta$ are taken when the
gradient of \eqref{eq:ch26-pushforward} is formed. In effect the objective evaluates the map
on a set that contains the trajectory, so the one-step accuracy it buys is accuracy in the
region the bound needs. That is the whole of the argument, and it is a plausibility
argument, not the theorem Part~I asked for: nothing in
\chapref{ch:part1-open-theoretical-problems-and-mathematical} or here relates the weights
$\lambda_k$ and the horizon $K$ to the Lipschitz product \eqref{eq:ch09-lipschitz}, or says
what condition on $(\lambda_k, K)$ makes training to horizon $K$ imply controlled error to
horizon $\alpha K$. \index{pushforward regularization!no convergence theorem}

The cost of \eqref{eq:ch26-pushforward} grows with the horizon it is asked to control, and
that growth is what decides how it is implemented. A term at index $k$ requires $k$ forward
applications, so one gradient step over the
full sum costs $\mathcal{O}(K)$ transforms and, if differentiated exactly, $\mathcal{O}(K)$
stored activations; per \eqref{eq:ch8-costs} each application is itself
$\Theta(d_v N\log N)$ plus the per-mode contraction. The standard compromise is truncated
backpropagation: unroll $K$ steps, detach the intermediate states, and let only the final
segment of the chain carry gradients. The consequence deserves naming, since it is a
property of the recipe and not a bug: the network is trained to be accurate \emph{at} the
states it reaches, but not to be accurate \emph{by the route} by which it reaches them, so
the composed map's stability is only weakly supervised. Full-chain gradients over long
horizons are unusable for a second reason beyond cost, since the backward pass multiplies by
the per-layer constants of \eqref{eq:ch09-lipschitz} once per unrolled step, and the same
$\Lambda_\theta$ that makes a rollout divergent makes the gradient explode or vanish.

Pushforward fails in three recognizable ways. If the target dynamics are chaotic --- the
leading Lyapunov exponents of the time-$\Delta t$ map positive, which is the hypothesis the
attractor discussion around
\eqref{eq:ch09-absorbing} records for two-dimensional Navier--Stokes flow --- then beyond the
predictability horizon the true state
$u(k\Delta t)$
is effectively uncorrelated with anything reachable from $u_0$, and minimizing a squared
pathwise error against an unpredictable target returns the conditional mean, which for a
turbulent vorticity field is a damped, near-laminar state. The rollout then stays bounded,
smooths, and is wrong in the respect a statistical benchmark is built to detect, so the
accounting is the weak one of
\eqref{eq:ch09-weak}: the bias of the learned invariant measure, not $e_K$. So $\lambda_k$
must decay in $k$, or $K$ must be capped near the horizon the physics allows --- and the
choice of where is a claim about the dynamics; no tuning of the optimizer settles it. Second, with weights
concentrated on large $k$ the objective is dominated by terms whose gradients pass through
many applications of the network, which is where \eqref{eq:ch09-lipschitz} predicts the
worst conditioning; uniform $\lambda_k$ is therefore the aggressive setting, and a decaying
$\lambda_k$ the conservative one. Third, pushforward trains on the union of a finite
set of model trajectories, so the
input distribution it optimizes is narrower than the one it is drawn from, and the
generalization the ensemble normalization of
\chapref{ch:part2-the-practitioner-s-paradigm-infinite-dimen} was supposed to provide is
partly forfeited: the normalizer's mean and standard deviation fields describe the solver's
ensemble, and a drifting model trajectory is measured against a reference the model is not
drawing from. \index{truncated backpropagation!through the rollout}
\index{chaotic dynamics!predictability horizon for a surrogate}
\index{normalization!drift of a rollout against the training ensemble}

% -----------------------------------------------------------------------------
\section{Failure Modes: A Field Guide}
\label{sec:ch26-failures}
% -----------------------------------------------------------------------------

What follows is arranged as symptom, then mechanism, then remedy, in the order a debugging
session proceeds: the band, the arithmetic, the horizon, the grid, the loader, and last the
failure that leaves every curve falling. Each entry
states a mechanism established above, not a measurement. \index{failure modes!field guide}

Validation error sits at the same level for mode budgets of eight, sixteen, and twenty-four
while the training error keeps falling. Read the run through
\eqref{eq:ch05-error-decomposition}: either the truncation term is no longer binding, or the
schedule stopped the band edge from being fitted at all, and the two look identical on a loss
curve. The discriminator is the slope Figure~\ref{fig:spectral_decay}(b) computes --- doubling
the retained band should divide the truncation contribution by roughly $2^{\,s-1/2}$ for a
target of regularity $s$, that is by about $1.4$, $2.8$, and $5.7$ at $s = 1, 2, 3$ --- so a
mode-budget sweep must be run with \emph{steps per run held proportional to the band}, as
Section~\ref{sec:ch26-schedule} argues, not with the schedule of
Listing~\ref{lst:ch26-onecycle} fixed. If widening the band still does
nothing once the steps are scaled, the binding term is capacity or one of the statistical and
discretization contributions separated in \cite{subedi2024controlling}, and no further mode
budget will move it.

The loss becomes \texttt{NaN} or the gradients become infinite a few hundred steps in, on a
GPU and with AMP enabled, and the same run is uneventful on the host central processing
unit (CPU) in single precision. The
mechanism is the array shift of Section~\ref{sec:ch26-amp}: a coefficient that leaves
binary16's representable range becomes an infinity, the subsequent contraction propagates it,
and the first non-finite entry is in a frequency intermediate rather than in an output field.
Locate the first tensor containing a non-finite value to confirm which layer, then act on the
range and not on the model: switch the autocast dtype to \texttt{bfloat16}, keep the transform
and the complex product at \texttt{float32}, and if binary16 must be used, add gradient
scaling so that the backward pass does not flush to zero. An infinite gradient that appears
only after resolution refinement is the same mechanism with the grid as the trigger.

Single-step error is excellent and the rollout dies around step twenty. Compute
\eqref{eq:ch09-lipschitz} at the checkpoint; the right-hand side is a product of quantities
readable from the parameters, and Proposition~\ref{prop:ch09-lip} makes it an upper bound on
$\Lambda_\theta$ for the band-limited realization. If it exceeds one and grows through
training, the objective is satisfied and the dynamics are not, because nothing in the loss
penalizes the product. Remedies in increasing order of cost: reduce the peak rate or deepen
the anneal so fewer expanding directions are learned late, add pushforward terms at modest
$K$ with geometrically decaying $\lambda_k$, and check whether the model has not blown up but
collapsed to a smooth state, in which case the failure is the conditional-mean one of
Section~\ref{sec:ch26-rollout} and the run should be judged by \eqref{eq:ch09-weak} rather
than by any pathwise curve. \index{Lipschitz product!diagnosing a divergent rollout}

Zero-shot evaluation on a refined grid returns a plausible but visibly smoothed version of the
truth. The retained band is a fixed set of array positions, and on a periodic domain of
unchanged length those positions are fixed physical wavenumbers: refining the grid raises the
Nyquist frequency of the array the truth lives on, while the spectral branch still addresses
only the modes \eqref{eq:ch05-modeset} gives it and the activations add nothing above them
that was ever fitted. The output is
therefore a low-pass copy of the fine-grid truth --- the model is not broken, it is being
asked for content it never represented. The
super-resolution results reported for operator models hold under a condition that is easy to
forget --- the inputs are band-limited inside the retained set, so what the experiment
demonstrates is resolution transfer, not recovery of detail absent from the coarse sampling,
and \eqref{eq:limit_commute} should be read with that premise attached. If the application
genuinely needs the content above $k_{\max}$, the remedy is to widen the retained band and
retrain, or to train at more than one resolution; interpolating the coarse prediction cannot
supply what the band excludes.
\index{resolution transfer!versus recovery of lost detail}

The reported error floor moves when the batch size or the number of GPUs changes, with
nothing else altered. Two mechanisms are entangled and should be checked in this order. The
loss \eqref{eq:rel_l2_loss} normalizes per sample before averaging, so low-amplitude members
of a batch receive large weight in the gradient; the small guard added to the denominator in
the implementation prevents a division blow-up but does not prevent that weighting, and a
reshuffled batch changes the mixture. Then the schedule, whose one invariant is that
\texttt{total\_steps} must equal the number of \texttt{scheduler.step()} calls the run will
actually make. Doubling the batch size halves \texttt{len(train\_loader)}, and the two ways to
react to that fail differently: recompute \texttt{total\_steps} and the cosine keeps its shape
while the run makes half as many updates of twice the data, leave it stale and the run stops
mid-descent with the rate still high and the anneal never completed. Only the second is a
mistake, and it is caught by logging the rate, not the loss: a run that ends at a learning
rate its author did not intend has not finished annealing. Either way a batch-size study must
fix the number of optimizer steps as well as the wall clock, or it measures the schedule and
not the batch.
\index{relative $L^2$ loss!per-sample weighting}

A sixth entry, because it is silent and permanent: training error falls, validation error
falls, and the network stops using the spectrum at all. Weight decay acts on
$R_{\phi,\ell}(k)$ toward zero without touching the phase, and a multiplier driven small
leaves the pointwise bypass $W_\ell$ as the only contentful part of the layer --- at which
point the model is a shallow feedforward map on the grid, accurate on the training resolution
and stripped of the global coupling that made it an operator. The diagnostic is direct: plot
$\sup_k \lVert R_{\phi,\ell}(k)\rVert$ against $\lVert W_\ell\rVert$ per layer per epoch, and
reduce the decay coefficient if the ratio trends downward through training.
\index{weight decay!suppression of the spectral branch}

\takeaway{What This Chapter Establishes}{Three floors bound a spectral surrogate before any
theorem does. The schedule, because gradient magnitude across the multiplier array falls
algebraically by \eqref{eq:ch26-mode-gradient}, and annealing suppresses precisely the band
edge with the least signal. The format, because an unscaled forward transform puts the
coefficient array's dynamic range beyond binary16 at the top for large grids and beyond it at
the bottom for steep spectra. And the horizon, because \eqref{eq:ch26-pushforward} is the only
term in the objective that sees the states a rollout actually visits, and it sees them without
bounding $\Lambda_\theta$. Each is checkable from parameters, tensor dtypes, or a loss curve,
which is what makes them worth checking before the physics is blamed.}

% -----------------------------------------------------------------------------
% Bridge to Chapter 16
% -----------------------------------------------------------------------------
\medskip
The three floors above are all properties of a training run, and a run ends. What has to
survive it is the parameter vector, exported into a runtime that no longer contains a
Python scheduler, a dtype policy, or an author who knows which half of the network was kept
in single precision and why. That is the subject of
\chapref{ch:part2-production-deployment-and-high-throughput-}: freezing a trained operator
into a traced graph, moving it off the training stack, and measuring what one application
costs against the classical solver it replaces. Two of this chapter's conclusions carry over
directly and inconveniently. An autocast region is Python-side state, so the archive records
the dtypes a trace produced and not the policy that produced them; the boundary cast this
chapter recommends must therefore be re-verified on the exported graph, where the gain from
reduced precision is a measurement made on that graph rather than a property carried with the
weights. And the zero-shot refinement that
chapter checks on its exported model --- a graph traced on a $64^2$ input and then applied to
$256^2$ --- is exactly the case in which Section~\ref{sec:ch26-failures} warns that a
correctly functioning operator returns a low-pass copy of the truth. Which of the failures
catalogued here become acceptance tests, and at what latency, is the next chapter's question.
